\scanpage{56}
\mspage{52}
% archive sheet number: 55
d'où
\[
\Sigma_{\unsure{\mathcal A|\mathscr A}}\simeq\Sigma_{\Sigma}\simeq\Sigma
\]
et on voit que ça respecte relations \unsure{d'ordres|d'ordre} et $C$.

Ceci dit, l'application canonique
\[
\mathcal A\longrightarrow\deleted{\unsure{\mathrm F}}\unsure{\mathrm{At}|\mathrm{Ft}}\,\mathrm{spat}(\Sigma_{\mathcal A})
\]
s'identifie, via $\Sigma_{\mathcal A}\simeq\Sigma$, en \unsure{composant}
\diagram{\begin{tikzcd}[ampersand replacement=\&,column sep=1.8em]
\mathcal A \arrow[r] \& \mathrm{Atspat}(\Sigma_{\mathcal A}) \arrow[r,"\sim"] \& \mathrm{Atspat}(\Sigma)=\mathcal A
\end{tikzcd}}
à l'application identique\unsure{,|.} Cette application canonique est donc non seulement injective, mais \uline{bijective} !

Revenons au cas d'un \unsure{Atelier|atelier} \deleted{spatial} quelconque $\mathcal A$. \deleted{\unsure{Il est}}\inserted{Posons $\Sigma=\Sigma_{\mathcal A}$.} Alors \deleted{$\mathcal A$ est \unsure{réalisé} comme sous-\unsure{Atelier|atelier} de $\mathrm{Atspat}(\Sigma)$}\inserted{on a un \uline{hom} \uline{\unsure{d'Ateliers|d'ateliers}}}
\[
\begin{aligned}
\mathcal A&\longrightarrow\mathrm{Atspat}(\Sigma)\\
X&\longmapsto\text{Dép}\,(\unsure{x|X})\qquad ,
\end{aligned}
\]
(même si \deleted{$\unsure{\Sigma|X}$} $\mathcal A$ non spatial). C'est un hom \uline{\unsure{fidèle|fidele}}, i.e.
\[
X\mathrel{|\unsure{o|\circ}|}Y\Longleftrightarrow\text{\unsure{Dep|Dép}}\,(x)\mathrel{|o|}\text{Dep}\,Y.
\]
Dire que $\mathcal A$ est spatial, signifie
